\CSBHeading{Dedekinds Originalbeweis}
Richard Dedekinds Manuskript \emph{Ähnliche (deutliche) Abbildung und
ähnliche Systeme} ist auf den 11.~Juli 1887 datiert. Es erschien 1932
in seinen \emph{Gesammelten mathematischen Werken}, Band~III,
S.~447--449
(\href{https://rcin.org.pl/Content/142326/PDF/WA35_176323_15883-3_Art17.pdf}{Originalscan}).
Seine Konstruktion lässt sich in heutiger Notation so beschreiben:

Sei \(\varphi\colon M\inj M\) und
\(\varphi[M]\subseteq N\subseteq M\). Setze \(U:=M\setminus N\).
Dedekind betrachtet die von \(U\) erzeugte Kette \(U_0\), also
die kleinste Teilmenge von \(M\), die \(U\) enthält und unter
\(\varphi\) abgeschlossen ist. Für sie gilt
\[
  U_0=U\cup\varphi[U_0].
\]
Daraus folgt die disjunkte Zerlegung
\[
  N=\varphi[U_0]\cup(M\setminus U_0).
\]
Die Abbildung
\[
  \psi(m):=\begin{cases}
    \varphi(m),&m\in U_0,\\
    m,&m\in M\setminus U_0
  \end{cases}
\]
bildet die beiden Teile von \(M\) bijektiv auf diese beiden Teile
von \(N\) ab. Damit ist \(\psi\colon M\bij N\).

\Needspace{14\baselineskip}
\CSBHeading{Vergleich mit der vorliegenden Fassung}
Setzt man
\[
  M=A,\qquad N=G[B],\qquad \varphi=G\circ F,
\]
so ist \(U=A\setminus G[B]\) und \(U_0=C\).
Die Bijektion \(\psi\colon A\bij G[B]\), gefolgt von der
Umkehrung der Bijektion \(G\colon B\bij G[B]\), ergibt genau
unsere Abbildung \(H\): auf \(C\) den Wert \(F(a)\), außerhalb
von \(C\) das eindeutige Urbild von \(a\) unter \(G\).

Die vorliegende Fassung führt somit dieselbe Beweisidee aus. Sie schreibt
die Schnittdefinition der erzeugten Menge und die Fixpunktgleichung
ausdrücklich hin und begründet die Bildgleichung sowie das Zusammenfügen
der beiden Bijektionen Schritt für Schritt.

\CSBHeading{Bernsteins Beweis in Borels Veröffentlichung}
Eine andere historische Darstellung steht bei Émile Borel,
\emph{Leçons sur la théorie des fonctions}, Paris: Gauthier-Villars et fils,
1898, S.~104--106
(\href{https://archive.org/details/leconstheoriefon00borerich/page/n117/mode/2up}{Originalscan ab S.~104}).
Borel schreibt den Beweis Felix Bernstein zu und weist darauf hin,
dass er die ihm übermittelte Fassung verändert hat, um transfinite
Zahlen zu vermeiden.

Die Schichtenkonstruktion lässt sich mit unseren Injektionen
\(F\colon A\inj B\) und \(G\colon B\inj A\) folgendermaßen
ausführen. Setze \(T:=G\circ F\) und definiere für \(n\geq0\)
\[
  A_0:=A,\qquad A_1:=G[B],\qquad A_{n+2}:=T[A_n].
\]
Wegen \(T[A]\subseteq G[B]\subseteq A\) und der Monotonie der
Bildbildung ist diese Folge absteigend.
Die \emph{Schichten} und der gemeinsame \emph{Rest} sind
\[
  D_n:=A_n\setminus A_{n+1},\qquad
  R:=\bigcap_{n\geq0}A_n.
\]
Jeder Punkt von \(A\) liegt entweder in allen \(A_n\), also in
\(R\), oder fällt an einer ersten Stelle aus der Folge heraus und
gehört damit zu genau einer Schicht \(D_n\). Die Schichten und
\(R\) zerlegen somit \(A\) in disjunkte Teile.

Da \(T\) injektiv ist, bildet es jede Schicht bijektiv auf die
übernächste ab:
\[
  T[D_n]=T[A_n]\setminus T[A_{n+1}]
        =A_{n+2}\setminus A_{n+3}=D_{n+2}.
\]
Die Idee ist nun, nur die geraden Schichten zu verschieben:
\(D_0\) nach \(D_2\), \(D_2\) nach \(D_4\) und so weiter.
Die ungeraden Schichten und der gesamte Rest \(R\) bleiben fest.
Dadurch wird genau \(D_0=A\setminus G[B]\) aus dem Bild entfernt.

\Needspace{12\baselineskip}
Fasse dazu die geraden Schichten zu
\(E:=\bigcup_{k\geq0}D_{2k}\) zusammen. Es gilt
\(T[E]=E\setminus D_0\). Die Abbildung
\[
  \psi(a):=\begin{cases}
    T(a),&a\in E,\\
    a,&a\in A\setminus E
  \end{cases}
\]
ist auf jedem der beiden Teile injektiv. Ihre Bildmengen
\(E\setminus D_0\) und \(A\setminus E\) sind disjunkt und
ergeben zusammen
\[
  (E\setminus D_0)\cup(A\setminus E)
    =A\setminus D_0=G[B].
\]
Damit ist \(\psi\colon A\bij G[B]\).
Wie bei Dedekind folgt anschließend die Umkehrung der Bijektion
\(G\colon B\bij G[B]\). Das ergibt \(H(a)=F(a)\) auf \(E\)
und das eindeutige \(G\)-Urbild von \(a\) außerhalb von \(E\).
Unter \(\psi\) bleibt \(R\) punktweise fest; im letzten Schritt
werden auch seine Punkte durch diese Umkehrung nach \(B\) übertragen.

\CSBHeading{Zusammenhang mit der Schnittkonstruktion}
Die Menge \(E\) ist genau die zuvor konstruierte Menge \(C\).
Sie enthält \(D_0=A\setminus G[B]\) und ist unter \(T\)
abgeschlossen. Umgekehrt muss jede Menge mit diesen beiden Eigenschaften
neben \(D_0\) auch \(T[D_0]=D_2\), dann \(D_4\) und induktiv
alle geraden Schichten enthalten. Daher ist \(E\) die kleinste
solche Menge, also \(E=C\). Die Schnittkonstruktion bestimmt somit
dieselbe Menge, ohne die Folge der Schichten vorauszusetzen.

\CSBHeading{Formale Fassung im Projekt}
Die nummerierten Aussagen stehen in Band~08; der Hauptsatz ist
\FormulaRefAuto{CantorBernsteinFixedInjections}[thm].
Die vollständigen formalen Ableitungen sind in den verknüpften
\CSBProofsLink{} ausgeführt.
